Applying
@@M0@@ gives the exact sequence:
@@M1@@
Since @@M2@@ by hypothesis (G), we obtain:
\sgaStableTarget{sga6.e08.eq.2-6-2}@@M3@@

Combining (2.6.1) and (2.6.2), we are reduced to checking that the diagram:
@@M4@@
is commutative, which is done locally without difficulty by returning to the definitions of the morphisms involved.

From the resulting exact sequence of complexes, one deduces the announced distinguished triangle in the derived category.

\sgaStableTarget{sga6.e08.cor.2-7}\textbf{Corollary 2.7.} \uline{With the notation and hypotheses of} 2.6, \uline{we have in} @@M5@@ \uline{the following equality}:
@@M6@@

\sgaStableTarget{sga6.e08.sec.3}\subsection*{3. Riemann--Roch theorem: statement}

\sgaStableTarget{sga6.e08.para.3-1}\textbf{3.1.} In section 3, @@M7@@ is a quasi-compact scheme having an ample sheaf.

Let @@M8@@ be a projective complete-intersection morphism (1.1); suppose in addition that the virtual relative dimension of @@M9@@ (1.9)
